\subsection{The cube condition for free massive fields}\label{ssec:fffree}

This subsection proves that the averaged cube condition of
Definition~\ref{def:avgcube} holds for the free massive scalar and Proca
fields on every torus, with a constant that depends only on the dimension and
on the mass in units of the cube, hence uniformly in the lattice spacing at
fixed physical size of the cube. The proof has three steps: a reduction of
the inequality to linear functions of the field, which is exact for Gaussian
measures; the identification of the Dirichlet form on linear functions with
the additive Schwarz operator of the cubes; and an explicit stable
decomposition of a field into fields supported in the cubes, the parabolic
profile of Proposition~\ref{prop:gauss}. Upper bounds show that the constant
vanishes with the mass, so that the mass term is what makes the condition
hold.

\emph{Setting.} In this subsection $\mu>0$ is a lattice mass and $d\ge1$,
$n\ge2$ and $L\ge n+1$ are integers; $\mathbb T=(\Z/L\Z)^d$ is the torus,
$e_1,\dots,e_d$ its unit vectors, $E=\mathbb T\times\{1,\dots,d\}$ its set
of links, the link $(y,i)$ running from $y$ to $y+e_i$, and $P=\mathbb T
\times\{(i,j):1\le i<j\le d\}$ its set of plaquettes. For $f\in\R^{\mathbb
T}$ let $\nabla_if(y)=f(y+e_i)-f(y)$, let $d_0f=(\nabla_if(y))_{(y,i)}\in
\R^E$, and for $A\in\R^E$ let $A_i(y)=A(y,i)$ and
\[
(d_1A)(y,i,j)=A(y,i)+A(y+e_i,j)-A(y+e_j,i)-A(y,j),
\]
so that $d_1d_0=0$; $\delta_0=d_0^{\mathsf T}$ and $\delta_1=d_1^{\mathsf T}$
are the adjoints for the standard inner products, and $-\Delta=\delta_0d_0$
is the lattice Laplacian, $\ip f{-\Delta f}=\sum_i\lVert\nabla_if\rVert^2=:
\lVert\nabla f\rVert^2$. The scalar field has precision
$Q_{\rm s}=-\Delta+\mu^2$ on $\R^{\mathbb T}$ and the Proca field has
precision $Q_{\rm P}=\delta_1d_1+\mu^2$ on $\R^E$; both are positive
definite. For $x\in\mathbb T$ let $V_x=x+\{0,\dots,n-1\}^d$, a set of $n^d$
distinct sites since $L\ge n$, and let $E_x$ be the set of links with base
point in $V_x$. Every site lies in exactly $n^d$ sets $V_x$ and every link
in exactly $n^d$ sets $E_x$. The parameter that matters is $s=n\mu$: with
$\mu=ma$ and $n=\ell/a$ it is $m\ell$, the mass times the physical size of
the cube. The tori of Definition~\ref{def:avgcube}, $d=4$ and $L$ even with
$L\ge2n+4$, are included.

For a positive definite $Q$ on $\R^N$ let $\gamma_Q$ be the centred Gaussian
law with covariance $C=Q^{-1}$, and for $V\subset\{1,\dots,N\}$ let
$\mathcal P_V$ be the conditional expectation on $L^2(\gamma_Q)$ given the
coordinates outside $V$; it is an orthogonal projection. For a family
$(V_x)_{x\in X}$ of coordinate sets write $\mathcal P_x=\mathcal P_{V_x}$
and
\begin{equation}\label{eq:ffvarpi}
\varpi(Q;(V_x))=\inf_f\frac{n^{-d}\sum_{x\in X}\lVert f-\mathcal P_xf
\rVert^2}{\Var_{\gamma_Q}(f)},
\end{equation}
the infimum over nonconstant $f\in L^2(\gamma_Q)$. Thus $\mathrm{BP}(n,
\varpi)$ of Definition~\ref{def:avgcube}, with $L^2(\gamma_Q)$ in place of
$\mathcal A$, holds exactly when $\varpi\le\varpi(Q;(V_x))$. For the scalar
field the family is $(V_x)_{x\in\mathbb T}$, for the Proca field it is
$(E_x)_{x\in\mathbb T}$.

\subsubsection{Reduction to linear functions}

Let $H\subset L^2(\gamma_Q)$ be the space of linear functions $\ell_h(\phi)
=\ip h\phi$, $h\in\R^N$; $\E[\ell_h\ell_{h'}]=\ip h{Ch'}$. For $m\ge0$ let
$\mathcal Q_m\subset L^2(\gamma_Q)$ be the space of polynomials of degree at
most $m$ in $\phi$, and let $W_m=\mathcal Q_m\cap\mathcal Q_{m-1}^\perp$
($\mathcal Q_{-1}=\{0\}$), the $m$-th Wiener chaos; $W_0$ is the constants
and $W_1=H$. Polynomials are dense in $L^2(\gamma_Q)$~\cite[Chapter~2]{Jan97},
so $L^2(\gamma_Q)=\bigoplus_{m\ge0}W_m$, an orthogonal sum, and $\Var(f)=
\sum_{m\ge1}\lVert f_m\rVert^2$ for $f=\sum_mf_m$, $f_m\in W_m$. A subspace
$G\subset\R^N$ is called compatible with $Q$ if $QG=G$; then also
$QG^\perp=G^\perp$, since $Q$ is symmetric and invertible, so $Q$ commutes
with the orthogonal projection $P_{G^\perp}$ onto $G^\perp$. A function $f\in
L^2(\gamma_Q)$ is $G$-invariant if it is a function of $P_{G^\perp}\phi$,
that is, measurable with respect to $\sigma(P_{G^\perp}\phi)$; for a function
defined at every $\phi$ this is the same as $f(\phi+w)=f(\phi)$ for all
$\phi$ and all $w\in G$ (take $w=-P_G\phi$), which is gauge invariance when
$G$ is the space of gauge directions. The $G$-invariant functions form a
closed subspace $L^2_G$, and $\ell_h$ is $G$-invariant exactly when $h\in
G^\perp$.

\begin{theorem}[the gap is attained on linear functions]\label{thm:ffchaos}
Let $Q$ be positive definite on $\R^N$, let $(V_x)_{x\in X}$ be a finite
family of coordinate sets, and let $G$ be a subspace compatible with $Q$
(for instance $G=\{0\}$). For $\lambda\ge0$ the following are equivalent:
\begin{enumerate}[label=\textup{(\roman*)},leftmargin=*]
\item $\sum_{x\in X}\lVert f-\mathcal P_xf\rVert^2\ge\lambda\Var(f)$ for
every $f\in L^2_G$;
\item $\sum_{x\in X}\lVert f-\mathcal P_xf\rVert^2\ge\lambda\Var(f)$ for
every $f\in L^2_G\cap H$, that is, for $f=\ell_h$ with $h\in G^\perp$.
\end{enumerate}
In particular $\varpi(Q;(V_x))$ is the infimum in~\eqref{eq:ffvarpi} over
linear $f$, and the same holds for the infimum over $L^2_G$.
\end{theorem}

The proof uses the second quantization of the conditional expectations,
which we set up in finite dimension. Let $\xi$ be a standard Gaussian vector
on $\R^N$ with law $\gamma$, and write $\ell_u(\xi)=\ip u\xi$. The spaces
$\mathcal Q_m$ and $W_m$ are defined for $\gamma$ as above; they are
invariant under orthogonal changes of variables, since $\mathcal Q_m$ is.
Let $h_j$ be the Hermite polynomials, $h_0=1$, $h_1(t)=t$, $h_{j+1}(t)=
th_j(t)-jh_{j-1}(t)$; they are orthogonal in $L^2(\R,e^{-t^2/2}dt/\sqrt{2
\pi})$ with $\lVert h_j\rVert^2=j!$, and $h_j$ has degree $j$ and leading
coefficient $1$. For an orthonormal basis $(e_i)$ of $\R^N$, the variables
$\xi'_i=\ip{e_i}\xi$ are independent standard Gaussians, and for a
multi-index $\alpha\in\mathbb N^N$ we put $H_\alpha(\xi')=\prod_ih_{\alpha_i}
(\xi'_i)$, $\alpha!=\prod_i\alpha_i!$ and $\lvert\alpha\rvert=\sum_i
\alpha_i$. On $(\R^N)^{\otimes m}$ we use the inner product in which the
tensors $e_{i_1}\otimes\dots\otimes e_{i_m}$ are orthonormal, the
symmetrization $\mathrm{Sym}=(m!)^{-1}\sum_\sigma\sigma$ over the
permutations of the factors, and $\mathrm{Sym}^m(\R^N)$, its range.

\begin{lemma}[the chaos as symmetric tensors]\label{lem:fffock}
\begin{enumerate}[label=\textup{(\roman*)},leftmargin=*]
\item Let $(e_i)$ be an orthonormal basis and let $m\ge0$. Then
$\{H_\alpha(\xi'):\lvert\alpha\rvert=m\}$ is an orthogonal basis of
$W_m$.
\item There is a linear map $J_m:(\R^N)^{\otimes m}\to W_m$ with
$J_m(u_1\otimes\dots\otimes u_m)=\Pi_{W_m}(\ell_{u_1}\cdots\ell_{u_m})$,
where $\Pi_{W_m}$ is the orthogonal projection onto $W_m$. It satisfies
$J_m=J_m\circ\mathrm{Sym}$, and $(m!)^{-1/2}J_m$ restricted to
$\mathrm{Sym}^m(\R^N)$ is an isometry onto $W_m$.
\item Let $K\subset\R^N$ be a subspace, $\pi$ the orthogonal projection
onto $K$, $\mathcal F_K$ the $\sigma$-algebra generated by $\ell_u$, $u\in
K$, and $\E_K=\E[\,\cdot\mid\mathcal F_K]$. Then $\E_KJ_m(t)=J_m(\pi^{\otimes
m}t)$ for every $t\in(\R^N)^{\otimes m}$. In particular $\E_K$ maps $W_m$
into itself and $\E_K=J_m\pi^{\otimes m}J_m^{-1}$ on $W_m$, where $J_m^{-1}$
denotes the inverse of $J_m$ restricted to $\mathrm{Sym}^m(\R^N)$.
\end{enumerate}
\end{lemma}

\begin{proof}
(i) $H_\alpha(\xi')$ is a polynomial of degree $\lvert\alpha\rvert$ in
$\xi'$, hence in $\xi$, and $\mathcal Q_m$ is also the span of the monomials
$\xi'^\beta$, $\lvert\beta\rvert\le m$. For $\lvert\beta\rvert<\lvert\alpha
\rvert$ some $\beta_i<\alpha_i$, and by independence $\E[H_\alpha(\xi')
\xi'^\beta]=\prod_i\E[h_{\alpha_i}(\xi'_i)\xi_i'^{\beta_i}]=0$, since
$h_{\alpha_i}$ is orthogonal to the polynomials of lower degree. So
$H_\alpha(\xi')\in\mathcal Q_{\lvert\alpha\rvert}\cap\mathcal Q_{\lvert
\alpha\rvert-1}^\perp=W_{\lvert\alpha\rvert}$, and the $H_\alpha$ are
mutually orthogonal, with $\lVert H_\alpha\rVert^2=\alpha!$. Since
$\xi'^\alpha-H_\alpha(\xi')$ has degree less than $\lvert\alpha\rvert$,
induction on the degree gives $\mathcal Q_m=\mathrm{span}\{H_\beta:\lvert
\beta\rvert\le m\}$; so $\{H_\beta:\lvert\beta\rvert\le m\}$ is an orthogonal
basis of $\mathcal Q_m$, and $\{H_\alpha:\lvert\alpha\rvert=m\}$ one of
$\mathcal Q_m\ominus\mathcal Q_{m-1}=W_m$.

(ii) The map $(u_1,\dots,u_m)\mapsto\Pi_{W_m}(\ell_{u_1}\cdots\ell_{u_m})$ is
multilinear and symmetric, so it defines a linear $J_m$ on $(\R^N)^{\otimes
m}$ with $J_m=J_m\circ\mathrm{Sym}$. Fix an orthonormal basis $(e_i)$. For
$i_1,\dots,i_m$ with multiplicity vector $\alpha$, $\ell_{e_{i_1}}\cdots
\ell_{e_{i_m}}=\xi'^\alpha$ and $\xi'^\alpha-H_\alpha(\xi')\in\mathcal
Q_{m-1}$, so $J_m(e_{i_1}\otimes\dots\otimes e_{i_m})=H_\alpha(\xi')$ by
(i). The tensors $S_\alpha=(m!/\alpha!)^{1/2}\mathrm{Sym}(e_1^{\otimes
\alpha_1}\otimes\dots\otimes e_N^{\otimes\alpha_N})$, $\lvert\alpha\rvert=
m$, form an orthonormal basis of $\mathrm{Sym}^m(\R^N)$: $\mathrm{Sym}
(e^{\otimes\alpha})$ is $\alpha!/m!$ times the sum of the $m!/\alpha!$
distinct basis tensors with multiplicity vector $\alpha$, so these
symmetrizations are orthogonal for distinct $\alpha$, have squared norm
$\alpha!/m!$, and span $\mathrm{Sym}^m(\R^N)$. Then $J_m(S_\alpha)=(m!/
\alpha!)^{1/2}H_\alpha(\xi')$, and these are orthogonal with squared norm
$m!$ and span $W_m$ by (i).

(iii) Both sides are linear in $t$, so it suffices to take $t=e'_{i_1}
\otimes\dots\otimes e'_{i_m}$ for an orthonormal basis $(e'_i)$ with
$e'_1,\dots,e'_k$ spanning $K$; such a basis exists, and (ii) holds for it
because $J_m$ was defined without reference to a basis. The $\sigma$-algebra
$\mathcal F_K$ is generated by $\xi'_1,\dots,\xi'_k$, and the $\xi'_i$ are
independent, so for the multiplicity vector $\alpha$ of $(i_1,\dots,i_m)$
\[
\E_KH_\alpha(\xi')=\prod_{i\le k}h_{\alpha_i}(\xi'_i)\prod_{i>k}\E[h_{
\alpha_i}(\xi'_i)],
\]
which is $H_\alpha(\xi')$ if $\alpha_i=0$ for all $i>k$ and $0$ otherwise,
since $\E h_j(\xi'_i)=0$ for $j\ge1$. On the other side $\pi e'_i=e'_i$ for
$i\le k$ and $\pi e'_i=0$ for $i>k$, so $J_m(\pi e'_{i_1}\otimes\dots\otimes
\pi e'_{i_m})$ is $H_\alpha(\xi')$ if all $i_l\le k$ and $0$ otherwise. The
two agree. The last statement follows from (ii).
\end{proof}

\begin{proof}[Proof of Theorem~\ref{thm:ffchaos}]
(i) implies (ii) since $L^2_G\cap H\subset L^2_G$. Assume (ii). Let $\xi=
C^{-1/2}\phi$, a standard Gaussian vector, so that $\ell_h(\phi)=\ip{C^{1/2}
h}{\xi}$, and $f\mapsto f(C^{1/2}\,\cdot)$ is an isometry of $L^2(\gamma_Q)$
onto $L^2(\gamma)$ that maps $\mathcal Q_m$, $W_m$ and $H$ onto the
corresponding spaces for $\xi$. For $x\in X$ let $W=V_x^c$. The
$\sigma$-algebra generated by $\phi_y$, $y\in W$, is generated by $\ip{C^{1/2}
\delta_y}\xi$, $y\in W$, that is, it is $\mathcal F_{K_x}$ with $K_x=C^{1/2}
\R^W$, where $\R^W$ is the set of vectors supported in $W$. So $\mathcal P_x
=\E_{K_x}$, and by Lemma~\ref{lem:fffock}(iii) it preserves each $W_m$ and
acts on it as $J_m\pi_x^{\otimes m}J_m^{-1}$, with $\pi_x$ the orthogonal
projection onto $K_x$.

Let $B=\sum_x(I-\pi_x)$ on $\R^N$. For $u\in\R^N$, $\ell_u(\xi)=J_1(u)$ and
$\lVert\ell_u-\mathcal P_x\ell_u\rVert^2=\lVert(I-\pi_x)u\rVert^2$, so (ii)
for $f=\ell_h$, $h\in G^\perp$, reads $\ip u{Bu}\ge\lambda\lvert u\rvert^2$
for $u\in T:=C^{1/2}G^\perp$. As $Q$ commutes with $P_{G^\perp}$, so does $C^{1/2}$ (a function of
$Q$), and $T=G^\perp$.

Fix $m\ge1$ and write $\pi_{(j)}$ for the operator on $(\R^N)^{\otimes m}$
that acts as $\pi$ on the $j$-th factor and as the identity on the others.
The $\pi_{(j)}$, $j=1,\dots,m$, are commuting orthogonal projections and
$\pi^{\otimes m}=\prod_j\pi_{(j)}$ is the orthogonal projection onto the
intersection of their ranges, so $\pi^{\otimes m}\le\pi_{(j)}$ for each $j$,
and averaging over $j$,
\begin{equation}\label{eq:fftensor}
I-\pi^{\otimes m}\ge\frac1m\sum_{j=1}^m(I-\pi_{(j)}).
\end{equation}
Applying this to each $\pi_x$ and summing over $x$ gives $\sum_x(I-\pi_x^{
\otimes m})\ge\frac1m\sum_jB_{(j)}$. For $t\in T^{\otimes m}$ and each $j$,
$\ip t{B_{(j)}t}\ge\lambda\lvert t\rvert^2$: writing $t=\sum_{a}t_a\otimes
f_a$ with $(f_a)$ an orthonormal basis of $T$ in the $j$-th slot and $t_a$
tensors in the other slots, $\ip t{B_{(j)}t}=\sum_{a,b}\ip{t_a}{t_b}\ip{f_a}
{Bf_b}=\operatorname{tr}(\Gamma B_T)$, where $\Gamma_{ab}=\ip{t_a}{t_b}$ and
$(B_T)_{ab}=\ip{f_a}{Bf_b}$ are positive semidefinite and $B_T\ge\lambda$ by
the hypothesis on $T$; so $\operatorname{tr}(\Gamma B_T)\ge\lambda
\operatorname{tr}\Gamma=\lambda\lvert t\rvert^2$. Hence
\begin{equation}\label{eq:ffchaosm}
\sum_x\ip t{(I-\pi_x^{\otimes m})t}\ge\lambda\lvert t\rvert^2\qquad(t\in
T^{\otimes m}).
\end{equation}

Now let $f\in L^2_G$ with chaos components $f_m\in W_m$. Choose an
orthonormal basis $(e'_i)$ of $\R^N$ with $e'_1,\dots,e'_k$ spanning $T$.
Since $f$ is a function of $P_{G^\perp}\phi$ and $P_{G^\perp}C^{1/2}\xi=
C^{1/2}P_T\xi$, the function $f(C^{1/2}\xi)$ is a function of $P_T\xi$, that
is, of $\xi'_1,\dots,\xi'_k$ alone. Its Hermite coefficients
$\E[f(C^{1/2}\xi)H_\alpha(\xi')]$ vanish unless $\alpha_i=0$ for all $i>k$,
by independence and $\E h_j(\xi'_i)=0$ for $j\ge1$; since the normalized
$H_\alpha(\xi')$ form an orthonormal basis of $L^2(\gamma)$ by
Lemma~\ref{lem:fffock}(i) and the density of polynomials, $f(C^{1/2}\xi)$
is the sum of its expansion in them. So each $f_m$ lies in the span of $H_\alpha(\xi')$ with $\lvert\alpha
\rvert=m$ and $\alpha$ supported in $\{1,\dots,k\}$, which by
Lemma~\ref{lem:fffock}(ii) is $J_m(\mathrm{Sym}^m(T))$. Write $f_m=
(m!)^{-1/2}J_m(t_m)$ with $t_m\in\mathrm{Sym}^m(T)\subset T^{\otimes m}$.
Since $\mathcal P_x$ preserves the orthogonal decomposition into chaoses,
$\lVert f-\mathcal P_xf\rVert^2=\sum_m\lVert f_m-\mathcal P_xf_m\rVert^2$,
and $\lVert f_m-\mathcal P_xf_m\rVert^2=\lvert t_m-\pi_x^{\otimes m}t_m
\rvert^2=\ip{t_m}{(I-\pi_x^{\otimes m})t_m}$ by Lemma~\ref{lem:fffock}, using
that $\pi_x^{\otimes m}$ preserves $\mathrm{Sym}^m(\R^N)$ and that $I-\pi_x^{
\otimes m}$ is an orthogonal projection. By~\eqref{eq:ffchaosm},
\begin{align*}
\sum_x\lVert f-\mathcal P_xf\rVert^2&=\sum_{m\ge1}\sum_x\ip{t_m}{(I-\pi_x^{
\otimes m})t_m}\\
&\ge\lambda\sum_{m\ge1}\lvert t_m\rvert^2=\lambda\sum_{m\ge1}
\lVert f_m\rVert^2=\lambda\Var(f),
\end{align*}
which is (i). The statements about $\varpi$ follow with $\lambda=n^d\varpi$.
\end{proof}

\begin{remark}
The theorem is the analogue, for the sum of the projections $\mathcal P_x$,
of a result of Amit~\cite{Ami91} on their product, the Gibbs sampler with a
Gaussian target, whose rate is likewise determined by the first chaos. Nothing beyond the
finite torus is needed here: the field has finitely many coordinates and
there are finitely many cubes. The conditional expectation $\mathcal P_x$
need not map $L^2_G$ into itself; the proof only uses that it preserves the
chaos decomposition.
\end{remark}

\subsubsection{Linear functions and the additive Schwarz operator}

For $V\subset\{1,\dots,N\}$ let $R_V:\R^N\to\R^V$ be the restriction and
$Q_{VV}=R_VQR_V^{\mathsf T}$, a positive definite principal submatrix. In
the $Q$-inner product $\ip vw_Q=\ip v{Qw}$ the $Q$-orthogonal projection onto
$R_V^{\mathsf T}\R^V$, the fields supported in $V$, is $P_V^Q=R_V^{\mathsf T}
Q_{VV}^{-1}R_VQ$. Put
\begin{equation}\label{eq:ffschwarz}
S=\sum_{x\in X}R_{V_x}^{\mathsf T}Q_{V_xV_x}^{-1}R_{V_x},\qquad
P_{\rm ad}=\sum_{x\in X}P^Q_{V_x}=SQ,
\end{equation}
the additive Schwarz operator of the family~\cite{Lio88,TW05}.

\begin{lemma}[linear functions]\label{lem:fflinear}
For $h\in\R^N$ and $V\subset\{1,\dots,N\}$,
$\lVert\ell_h-\mathcal P_V\ell_h\rVert^2=\ip{h_V}{Q_{VV}^{-1}h_V}$ with
$h_V=R_Vh$. Hence $\sum_x\lVert\ell_h-\mathcal P_x\ell_h\rVert^2=\ip h{Sh}$,
$\Var(\ell_h)=\ip h{Q^{-1}h}$, and $n^d\varpi(Q;(V_x))$ is the smallest
eigenvalue of $Q^{1/2}SQ^{1/2}$, which is unitarily equivalent to $P_{\rm
ad}$ on $(\R^N,\ip\cdot\cdot_Q)$. If $G$ is compatible with $Q$, the infimum
in~\eqref{eq:ffvarpi} over $L^2_G$ is $n^{-d}$ times the smallest eigenvalue
of the compression of $Q^{1/2}SQ^{1/2}$ to $G^\perp$.
\end{lemma}

\begin{proof}
In the notation of the proof of Theorem~\ref{thm:ffchaos}, with $W=V^c$,
$\lVert\ell_h-\mathcal P_V\ell_h\rVert^2=\lvert(I-\pi)u\rvert^2$ with $u=
C^{1/2}h$ and $\pi$ the projection onto $K=C^{1/2}\R^W$. A vector $v$ is
orthogonal to $K$ exactly when $C^{1/2}v\perp\R^W$, that is $C^{1/2}v\in R_V^{
\mathsf T}\R^V$, so $K^\perp=Q^{1/2}R_V^{\mathsf T}\R^V$, and the orthogonal
projection onto $K^\perp$ is $I-\pi=Q^{1/2}R_V^{\mathsf T}Q_{VV}^{-1}R_V
Q^{1/2}$ (it is symmetric, idempotent, has range in $K^\perp$ and fixes
$Q^{1/2}R_V^{\mathsf T}w$). So $\lvert(I-\pi)u\rvert^2=\ip u{(I-\pi)u}=
\ip{Q^{1/2}u}{R_V^{\mathsf T}Q_{VV}^{-1}R_VQ^{1/2}u}=\ip{h_V}{Q_{VV}^{-1}
h_V}$, as $Q^{1/2}u=h$. Summing over $x$ gives $\ip h{Sh}$; the variance is
$\ip h{Ch}$. By Theorem~\ref{thm:ffchaos} with $G=\{0\}$, $n^d\varpi$ is the
infimum of $\ip h{Sh}/\ip h{Q^{-1}h}$, which is the smallest eigenvalue of
$Q^{1/2}SQ^{1/2}$ (put $h=Q^{1/2}u$); and $Q^{1/2}$ is a unitary map from
$(\R^N,\ip\cdot\cdot_Q)$ to $(\R^N,\ip\cdot\cdot)$ with $Q^{1/2}P_{\rm ad}
Q^{-1/2}=Q^{1/2}SQ^{1/2}$. For $G$ compatible with $Q$, $Q^{1/2}G^\perp=
G^\perp$, and the infimum over $h\in G^\perp$ is the smallest eigenvalue of
the compression to $G^\perp$.
\end{proof}

\begin{lemma}[stable decompositions, P.-L. Lions]\label{lem:fflions}
Let $T\subset\R^N$ be a subspace with $QT=T$ (for instance $T=\R^N$) and
$C_0>0$. Suppose that every $v\in T$ can be written as $v=\sum_{x\in X}v_x$
with $v_x$ supported in $V_x$ and $\sum_x\ip{v_x}{Qv_x}\le C_0\ip v{Qv}$.
Then $\ip h{Sh}\ge C_0^{-1}\ip h{Q^{-1}h}$ for every $h\in T$, and the
infimum in~\eqref{eq:ffvarpi} over $L^2_{T^\perp}$, in particular
$\varpi(Q;(V_x))$ when $T=\R^N$, is at least $1/(C_0n^d)$.
\end{lemma}

\begin{proof}
Let $h\in T$ and $v=Q^{-1}h\in T$. Since $v-P^Q_{V_x}v$ is $Q$-orthogonal to
the fields supported in $V_x$, $\ip{v_x}{Qv}=\ip{v_x}{QP^Q_{V_x}v}$, and
$\lVert P^Q_{V_x}v\rVert_Q^2=\ip v{QP^Q_{V_x}v}=\ip h{R_{V_x}^{\mathsf T}
Q_{V_xV_x}^{-1}R_{V_x}h}$. By the Cauchy--Schwarz inequality
\begin{align*}
\ip v{Qv}&=\sum_x\ip{v_x}{QP^Q_{V_x}v}\le\Bigl(\sum_x\lVert v_x\rVert_Q^2
\Bigr)^{1/2}\Bigl(\sum_x\lVert P^Q_{V_x}v\rVert_Q^2\Bigr)^{1/2}\\
&\le C_0^{1/2}\ip v{Qv}^{1/2}\ip h{Sh}^{1/2},
\end{align*}
so $\ip h{Q^{-1}h}=\ip v{Qv}\le C_0\ip h{Sh}$. The last claim follows from
Lemma~\ref{lem:fflinear}, since $T^\perp$ is compatible with $Q$ and $(T^\perp)^\perp=T$.
\end{proof}

\subsubsection{The partition of unity}

A profile is a function $\tau:\Z\to[0,\infty)$ supported in $\{0,\dots,n-1
\}$ with $\sum_s\tau(s)=1$; it is regarded as a function on $\Z/L\Z$, which
is possible since $L\ge n$. Put
\begin{gather*}
S_2=\sum_s\tau(s)^2,\qquad S_{11}=\sum_s\tau(s)\tau(s+1),\\
D_2=\sum_s(\tau(s+1)-\tau(s))^2=2(S_2-S_{11}),
\end{gather*}
the sums over $s\in\Z$, and $\alpha=nS_2$, $\beta=n^3D_2$. For $x\in\mathbb
T$ let $\chi_x(y)=\prod_{i=1}^d\tau(y_i-x_i)$. Then $\chi_x$ is supported in
$V_x$, $\sum_x\chi_x\equiv1$, and because $L\ge n+1$ the residues $-1,0,
\dots,n-1$ are distinct in $\Z/L\Z$, so that for all $y$ and $i\ne j$
\begin{equation}\label{eq:ffchisums}
\begin{gathered}
\sum_x\chi_x(y)^2=S_2^d,\qquad\sum_x\chi_x(y)\chi_x(y+e_i)=S_2^{d-1}S_{11},\\
\sum_x\chi_x(y+e_i)\chi_x(y+e_j)=S_2^{d-2}S_{11}^2 ,
\end{gathered}
\end{equation}
each sum over $x\in\mathbb T$ factorizing over the coordinates.

\begin{lemma}[the scalar identity]\label{lem:ffpartition}
For every profile and every $v\in\R^{\mathbb T}$,
\begin{align*}
\sum_{x\in\mathbb T}\ip{\chi_xv}{Q_{\rm s}\chi_xv}&=S_2^{d-1}\bigl[S_{11}
\lVert\nabla v\rVert^2+(S_2\mu^2+dD_2)\lVert v\rVert^2\bigr]\\
&\le S_2^{d-1}\Bigl(S_2+\frac{dD_2}{\mu^2}\Bigr)\ip v{Q_{\rm s}v}.
\end{align*}
\end{lemma}

\begin{proof}
For a link $(y,i)$, by~\eqref{eq:ffchisums},
\begin{multline*}
\sum_x\bigl(\nabla_i(\chi_xv)(y)\bigr)^2=\sum_x\bigl[\chi_x(y+e_i)^2v(y+e_i)^2
+\chi_x(y)^2v(y)^2\\
-2\chi_x(y)\chi_x(y+e_i)v(y)v(y+e_i)\bigr]\\
=S_2^d\bigl(v(y+e_i)^2+v(y)^2\bigr)-2S_2^{d-1}S_{11}v(y)v(y+e_i).
\end{multline*}
Summing over $y$, and using $2\sum_yv(y)v(y+e_i)=2\lVert v\rVert^2-\lVert
\nabla_iv\rVert^2$,
\begin{align*}
\sum_x\lVert\nabla_i(\chi_xv)\rVert^2&=2S_2^d\lVert v\rVert^2-S_2^{d-1}S_{11}
\bigl(2\lVert v\rVert^2-\lVert\nabla_iv\rVert^2\bigr)\\
&=S_2^{d-1}\bigl[S_{11}\lVert\nabla_iv\rVert^2+D_2\lVert v\rVert^2\bigr].
\end{align*}
Summing over $i$ and adding $\mu^2\sum_x\lVert\chi_xv\rVert^2=\mu^2S_2^d
\lVert v\rVert^2$ gives the identity. The inequality follows from $S_{11}\le
S_2$ (Cauchy--Schwarz) and $\ip v{Q_{\rm s}v}=\lVert\nabla v\rVert^2+\mu^2
\lVert v\rVert^2$.
\end{proof}

\begin{lemma}[the parabolic profile]\label{lem:ffparabola}
Let $\tau(s)=6(s+1)(n-s)/(n(n+1)(n+2))$ for $0\le s\le n-1$ and $\tau(s)=0$
otherwise. Then $\tau$ is a profile and
\begin{gather*}
\alpha_n:=nS_2=\frac{6(n^2+2n+2)}{5(n+1)(n+2)}\le\frac65,\\
\beta_n:=n^3D_2=\frac{12n^2}{(n+1)(n+2)}\le12,\\
\frac{\beta_n}{\alpha_n}=10\rho_n,\qquad\rho_n:=\frac{n^2}{n^2+2n+2}<1 .
\end{gather*}
\end{lemma}

\begin{proof}
With $u=s+1$ and $m=n+1$, the formulas for $\sum_{u=1}^nu^k$, $k\le4$, give
$\sum_{u=1}^nu(m-u)=n(n+1)(n+2)/6$ and $\sum_{u=1}^nu^2(m-u)^2=(m^5-m)/30=
n(n+1)(n+2)(n^2+2n+2)/30$, and the increments of $u\mapsto u(m-u)$ on $0\le
u\le n$ are $m-1-2u=n-2u$, with $\sum_{u=0}^n(n-2u)^2=n(n+1)(n+2)/3$. So
$\sum\tau=1$, $S_2=6(n^2+2n+2)/(5n(n+1)(n+2))$ and $D_2=12/(n(n+1)(n+2))$,
which give $\alpha_n$, $\beta_n$ and their ratio; $\alpha_n\le6/5$ since
$n^2+2n+2\le(n+1)(n+2)$, and $\beta_n\le12$ since $n^2\le(n+1)(n+2)$.
\end{proof}

With a constant profile $\tau=1/n$ one has $\alpha=1$ but $\beta=2n$, and the
bound below degenerates as $n\to\infty$; the smooth profile is what makes the
constant uniform in $n$.

\subsubsection{The scalar field}

\begin{theorem}[the cube condition for the free scalar field]\label{thm:ffscalar}
Let $d\ge1$, $n\ge2$, $L\ge n+1$, $\mu>0$ and $s=n\mu$. For the Gaussian
field on $\mathbb T$ with precision $Q_{\rm s}=-\Delta+\mu^2$ and the cubes
$V_x$, $x\in\mathbb T$,
\begin{multline*}
\varpi(Q_{\rm s};(V_x))\ge\frac1{\alpha_n^d+d\alpha_n^{d-1}\beta_ns^{-2}}
=\alpha_n^{-d}\frac{s^2}{s^2+10d\rho_n}\\
\ge\Bigl(\frac56\Bigr)^d\frac{s^2}{s^2+10d}=:c_d(s),
\end{multline*}
with $\alpha_n,\beta_n,\rho_n$ as in Lemma~\ref{lem:ffparabola}. That is,
$\mathrm{BP}(n,c_d(n\mu))$ holds for every real $f\in L^2(\gamma_{Q_{\rm
s}})$, and $c_d(s)\ge(5/6)^d(1+10d)^{-1}\min(1,s^2)$.
\end{theorem}

\begin{proof}
Let $\tau$ be the parabolic profile and $v_x=\chi_xv$. Then $v_x$ is
supported in $V_x$, $\sum_xv_x=v$, and by Lemma~\ref{lem:ffpartition}
$\sum_x\ip{v_x}{Q_{\rm s}v_x}\le C_0\ip v{Q_{\rm s}v}$ with $C_0=S_2^{d-1}
(S_2+dD_2/\mu^2)$. Lemma~\ref{lem:fflions} with $T=\R^{\mathbb T}$ gives
$\varpi\ge1/(C_0n^d)$, and $C_0n^d=(nS_2)^d+d(nS_2)^{d-1}n^3D_2/(n\mu)^2=
\alpha_n^d+d\alpha_n^{d-1}\beta_ns^{-2}$. The middle expression follows from
$\beta_n/\alpha_n=10\rho_n$, and the last inequality from $\alpha_n\le6/5$
and $\rho_n<1$. The claim about $\mathrm{BP}$ is the definition of
$\varpi$, and the final bound is $s^2/(s^2+10d)\ge\min(1,s^2)/(1+10d)$.
\end{proof}

For $d=4$ the bound is $\varpi\ge(5/6)^4s^2/(s^2+40)$, which is $0.0118$ at
$s=1$ and $0.0438$ at $s=2$. At $n=2$ the bound of the theorem, with the
$n$-dependent constants, is $s^2/(s^2+4d)$, and it coincides with the upper
bound of Proposition~\ref{prop:ffupper} below, so it is exact.

\subsubsection{The Proca field}

The gauge transformations of the Proca field are $A\mapsto A+d_0\psi$, $\psi
\in\R^{\mathbb T}$; let $G=d_0\R^{\mathbb T}$ be the space of gauge
directions, so that $G^\perp=\ker\delta_0$, the transverse fields. A
gauge-invariant function of $A$ is a $G$-invariant one, and a linear
function $\ell_h$ is gauge-invariant exactly when $\delta_0h=0$.

\begin{lemma}[the Hodge identity on the torus]\label{lem:ffhodge}
For every $A\in\R^E$, $\lVert d_1A\rVert^2+\lVert\delta_0A\rVert^2=\sum_{i=1}
^d\lVert\nabla A_i\rVert^2$. Consequently $Q_{\rm P}$ commutes with the
orthogonal projection onto $G$, so $G$ is compatible with $Q_{\rm P}$, and
for transverse $A$, $\ip A{Q_{\rm P}A}=\sum_i\lVert\nabla A_i\rVert^2+\mu^2
\lVert A\rVert^2$.
\end{lemma}

\begin{proof}
Put $\nabla_i^-f(y)=f(y)-f(y-e_i)$, so that $(\delta_0A)(y)=-\sum_i\nabla_i^-
A_i(y)$, and $(\nabla_i^-)^{\mathsf T}=-\nabla_i$ on $\R^{\mathbb T}$; all
these operators commute, being translations and their differences. Then
\begin{align*}
\lVert d_1A\rVert^2&=\sum_{i<j}\lVert\nabla_iA_j-\nabla_jA_i\rVert^2\\
&=\sum_{i\ne j}\lVert\nabla_jA_i\rVert^2-2\sum_{i<j}\ip{\nabla_iA_j}{\nabla_j
A_i},\\
\lVert\delta_0A\rVert^2&=\sum_i\lVert\nabla_i^-A_i\rVert^2+2\sum_{i<j}
\ip{\nabla_i^-A_i}{\nabla_j^-A_j}\\
&=\sum_i\lVert\nabla_iA_i\rVert^2+2\sum_{i<j}\ip{\nabla_jA_i}{\nabla_iA_j},
\end{align*}
since $\lVert\nabla_i^-f\rVert=\lVert\nabla_if\rVert$ and $\ip{\nabla_i^-
A_i}{\nabla_j^-A_j}=-\ip{A_i}{\nabla_i\nabla_j^-A_j}=-\ip{A_i}{\nabla_j^-
\nabla_iA_j}=\ip{\nabla_jA_i}{\nabla_iA_j}$. Adding the two lines gives
$\sum_{i,j}\lVert\nabla_jA_i\rVert^2$. Hence $\delta_1d_1=-\Delta\otimes I-
d_0\delta_0$. Now $\delta_1d_1$ vanishes on $G$ (as $d_1d_0=0$) and is
symmetric, so it maps $G^\perp$ into $G^\perp$; thus $Q_{\rm P}=\delta_1d_1+
\mu^2$ preserves $G$ and $G^\perp$, which is the commutation. For
transverse $A$ the identity reads $\lVert d_1A\rVert^2=\sum_i\lVert\nabla
A_i\rVert^2$.
\end{proof}

\begin{lemma}[the decomposition of a Proca field]\label{lem:ffprocadecomp}
Let $d\ge2$, let $\tau$ be a profile and $A_x(y,i)=\chi_x(y)A(y,i)$, so that
$A_x$ is supported in $E_x$ and $\sum_xA_x=A$. Then, with $\kappa_d=(5d-4)/2$,
\begin{align*}
\sum_{x\in\mathbb T}\ip{A_x}{Q_{\rm P}A_x}&\le S_2^{d-1}\bigl[S_2\lVert d_1A
\rVert^2+(S_2\mu^2+\kappa_dD_2)\lVert A\rVert^2\bigr]\\
&\le S_2^{d-1}\Bigl(S_2+\frac{\kappa_dD_2}{\mu^2}\Bigr)\ip A{Q_{\rm P}A},
\end{align*}
and if $\delta_0A=0$,
\begin{align*}
\sum_{x\in\mathbb T}\ip{A_x}{Q_{\rm P}A_x}&\le S_2^{d-1}\bigl[S_{11}\lVert d_1A
\rVert^2+(S_2\mu^2+dD_2)\lVert A\rVert^2\bigr]\\
&\le S_2^{d-1}\Bigl(S_2+\frac{dD_2}{\mu^2}\Bigr)\ip A{Q_{\rm P}A}.
\end{align*}
\end{lemma}

\begin{proof}
Fix a plaquette $p=(y,i,j)$ and put $a=A(y,i)-A(y,j)$, $b=A(y+e_i,j)$,
$c=A(y+e_j,i)$, so that $(d_1A)(p)=a+b-c$ and $(d_1A_x)(p)=\chi_x(y)a+\chi_x
(y+e_i)b-\chi_x(y+e_j)c$. With $M=S_2^d$, $M'=S_2^{d-1}S_{11}$ and $M''=
S_2^{d-2}S_{11}^2$, \eqref{eq:ffchisums} gives
\begin{align*}
\sum_x(d_1A_x)(p)^2&=M(a^2+b^2+c^2)+2M'(ab-ac)-2M''bc\\
&=M(d_1A)(p)^2-2(M-M')a(b-c)\\
&\quad+2(M-M'')bc.
\end{align*}
Here $M-M'=S_2^{d-1}(S_2-S_{11})=\frac12S_2^{d-1}D_2\ge0$ and $M-M''=S_2^{d-2}
(S_2-S_{11})(S_2+S_{11})\le S_2^{d-1}D_2$, using $S_{11}\le S_2$. Bound
$2\lvert a(b-c)\rvert\le a^2+(b-c)^2\le a^2+2b^2+2c^2$ and $2\lvert bc\rvert
\le b^2+c^2$, and sum over the plaquettes. For fixed $y$, $\sum_{i<j}(A_i(y)
-A_j(y))^2=d\sum_iA_i(y)^2-(\sum_iA_i(y))^2\le d\sum_iA_i(y)^2$, so $\sum_p
a^2\le d\lVert A\rVert^2$; and each link occurs as $b$ in $j-1$ plaquettes
and as $c$ in $d-i$ plaquettes, so $\sum_p(b^2+c^2)=(d-1)\lVert A\rVert^2$.
Therefore
\begin{align*}
\sum_x\lVert d_1A_x\rVert^2&\le M\lVert d_1A\rVert^2\\
&\quad+\bigl[(M-M')(3d-2)+(M-M'')(d-1)\bigr]\lVert A\rVert^2\\
&\le S_2^d\lVert d_1A\rVert^2+S_2^{d-1}D_2\Bigl[\frac{3d-2}2+d-1\Bigr]
\lVert A\rVert^2,
\end{align*}
and $\frac{3d-2}2+d-1=\kappa_d$. Adding $\mu^2\sum_x\lVert A_x\rVert^2=\mu^2
S_2^d\lVert A\rVert^2$ gives the first bound, and $\ip A{Q_{\rm P}A}=\lVert
d_1A\rVert^2+\mu^2\lVert A\rVert^2$ the second.

If $\delta_0A=0$, Lemma~\ref{lem:ffhodge} applied to $A_x$ and to $A$ gives
$\lVert d_1A_x\rVert^2\le\sum_i\lVert\nabla(\chi_xA_i)\rVert^2$ and $\sum_i
\lVert\nabla A_i\rVert^2=\lVert d_1A\rVert^2$. By the proof of
Lemma~\ref{lem:ffpartition},
\begin{align*}
\sum_x\sum_i\lVert\nabla(\chi_xA_i)\rVert^2
&=S_2^{d-1}\sum_i\bigl[S_{11}\lVert\nabla A_i\rVert^2+dD_2\lVert A_i\rVert^2\bigr]\\
&=S_2^{d-1}\bigl[S_{11}\lVert d_1A\rVert^2+dD_2\lVert A\rVert^2\bigr],
\end{align*}
and the mass term is as before.
\end{proof}

\begin{theorem}[the cube condition for the free Proca field]\label{thm:ffproca}
Let $d\ge2$, $n\ge2$, $L\ge n+1$, $\mu>0$ and $s=n\mu$. For the Gaussian
field on the links of $\mathbb T$ with precision $Q_{\rm P}=\delta_1d_1+\mu^2$
and the link cubes $E_x$, $x\in\mathbb T$:
\begin{enumerate}[label=\textup{(\roman*)},leftmargin=*]
\item for every real $f\in L^2(\gamma_{Q_{\rm P}})$, $\mathrm{BP}(n,\varpi)$
holds with
\[
\varpi=\alpha_n^{-d}\frac{s^2}{s^2+10\kappa_d\rho_n}\ge\Bigl(\frac56\Bigr)^d
\frac{s^2}{s^2+5(5d-4)},
\]
where $\kappa_d=(5d-4)/2$;
\item for every real gauge-invariant $f\in L^2(\gamma_{Q_{\rm P}})$,
$\mathrm{BP}(n,\varpi)$ holds with the constant of
Theorem~\ref{thm:ffscalar},
\[
\varpi=\alpha_n^{-d}\frac{s^2}{s^2+10d\rho_n}\ge c_d(s)=\Bigl(\frac56\Bigr)^d
\frac{s^2}{s^2+10d}.
\]
\end{enumerate}
In both cases the optimal constant is the infimum over linear functions of
the respective class, and the inequality in (ii) does not require $\mathcal
P_x$ to preserve gauge invariance.
\end{theorem}

\begin{proof}
(i) Lemma~\ref{lem:ffprocadecomp} with the parabolic profile gives a
decomposition of every $A$ with $C_0=S_2^{d-1}(S_2+\kappa_dD_2/\mu^2)$, and
Lemma~\ref{lem:fflions} with $T=\R^E$ gives $\varpi\ge1/(C_0n^d)$, $C_0n^d=
\alpha_n^d+\kappa_d\alpha_n^{d-1}\beta_ns^{-2}=\alpha_n^d(1+10\kappa_d\rho_n
s^{-2})$. (ii) $T=G^\perp$ satisfies $Q_{\rm P}T=T$ by
Lemma~\ref{lem:ffhodge}, and for $A\in T$ the second bound of
Lemma~\ref{lem:ffprocadecomp} gives $C_0=S_2^{d-1}(S_2+dD_2/\mu^2)$;
Lemma~\ref{lem:fflions} bounds the infimum over $L^2_G$, the
gauge-invariant functions, by $1/(C_0n^d)$, which is the constant of
Theorem~\ref{thm:ffscalar}. The statements about linear functions are
Theorem~\ref{thm:ffchaos} and Lemma~\ref{lem:fflinear}; the last remark was
made after Theorem~\ref{thm:ffchaos}.
\end{proof}

\subsubsection{Upper bounds}

\begin{proposition}[the mass is needed]\label{prop:ffupper}
Let $n\ge2$, $L\ge n+1$ and $s=n\mu$.
\begin{enumerate}[label=\textup{(\roman*)},leftmargin=*]
\item For the scalar field,
\begin{align*}
\varpi(Q_{\rm s};(V_x))&\le\frac{\mu^2}{n^d}\ip{\mathbf 1_V}{(-\Delta_D+\mu^2)
^{-1}\mathbf 1_V}\\
&\le\frac{s^2}{s^2+n^2\lambda_1(n)}\le\frac{s^2}{s^2+16d/9},
\end{align*}
where $-\Delta_D$ is the Dirichlet Laplacian of the cube $V=\{0,\dots,n-1
\}^d$ and $\lambda_1(n)=4d\sin^2(\pi/(2n+2))$ is its smallest eigenvalue.
\item For the Proca field with $d\ge2$, the infimum in~\eqref{eq:ffvarpi}
over gauge-invariant functions, hence also $\varpi(Q_{\rm P};(E_x))$, is at
most
\[
\frac{\mu^2}{n^d}\ip{h_V}{(Q_{\rm P})_{VV}^{-1}h_V}\le\frac{s^2}{1+s^2},
\]
where $V=E_0$ and $h_V$ is the indicator of the links of $V$ in the
direction $e_1$.
\end{enumerate}
In particular both constants are at most $s^2$, and tend to $0$ as
$n\mu\to0$.
\end{proposition}

\begin{proof}
(i) Take $h=\mathbf 1$, so that $Q_{\rm s}h=\mu^2h$ and $\Var(\ell_h)=
\ip h{Q_{\rm s}^{-1}h}=L^d/\mu^2$. For each $x$, Lemma~\ref{lem:fflinear} gives
\[
\lVert\ell_h-\mathcal P_x\ell_h\rVert^2=\ip{\mathbf 1_{V_x}}{(Q_{\rm s})
_{V_xV_x}^{-1}\mathbf 1_{V_x}},
\]
and $(Q_{\rm s})_{V_xV_x}=-\Delta_D+\mu^2$: two
sites of $V_x$ are neighbours on the torus exactly when they are neighbours
in $\Z^d$, since their coordinates differ by at most $n-1<L-1$. All $x$ give
the same value, so $\varpi\le n^{-d}L^d\ip{\mathbf 1_V}{(-\Delta_D+\mu^2)^{-1}
\mathbf 1_V}\mu^2/L^d$, the first bound. The eigenvalues of the Dirichlet
Laplacian of the path on $n$ sites are $2-2\cos(\pi k/(n+1))$, $1\le k\le n$,
and $-\Delta_D$ is the sum of $d$ commuting copies, so its smallest
eigenvalue is $\lambda_1(n)=d(2-2\cos(\pi/(n+1)))=4d\sin^2(\pi/(2n+2))$ and
$(-\Delta_D+\mu^2)^{-1}\le(\lambda_1+\mu^2)^{-1}$; with $\lVert\mathbf 1_V
\rVert^2=n^d$ this gives $\varpi\le\mu^2/(\lambda_1+\mu^2)=s^2/(s^2+n^2
\lambda_1)$. Finally $\sin t\ge2t/\pi$ on $[0,\pi/2]$ gives $n^2\lambda_1\ge
4dn^2/(n+1)^2\ge16d/9$ for $n\ge2$.

(ii) Let $h(y,i)=\delta_{i1}$, the constant field in the direction $e_1$.
Then $d_1h=0$ and $\delta_0h=0$, so $\ell_h$ is gauge-invariant, $Q_{\rm P}h=
\mu^2h$ and $\Var(\ell_h)=L^d/\mu^2$. Fix $x$, let $V=E_x$ and $h_V=R_Vh$,
with $\lVert h_V\rVert^2=n^d$, and write $I_i=\{x_i,\dots,x_i+n-1\}\subset
\Z/L\Z$. We claim that for every $w\in\R^E$ supported in $E_x$,
\begin{equation}\label{eq:ffclaim}
\ip{h_V}{w}^2\le n^{d+2}\lVert d_1w\rVert^2 .
\end{equation}
For $y_2\in\Z/L\Z$ and $\tilde y=(y_3,\dots,y_d)$ let
\begin{gather*}
W(y_2,\tilde y)=\sum_{y_1\in\Z/L\Z}w(y_1,y_2,\tilde y;1),\\
D(y_2,\tilde y)=\sum_{y_1\in\Z/L\Z}(d_1w)(y_1,y_2,\tilde y;1,2).
\end{gather*}
In the sum defining $D$ the terms $w(y+e_1,2)
-w(y,2)$ telescope over the period, so $D(y_2,\tilde y)=W(y_2,\tilde y)-W(y_2
+1,\tilde y)$. As $w$ is supported in $E_x$, $W(y_2,\tilde y)=0$ unless $y_2
\in I_2$ and $\tilde y\in I_3\times\dots\times I_d$; and $x_2-1\notin I_2$
since $L\ge n+1$. Hence $W(x_2+k,\tilde y)=-\sum_{j=0}^kD(x_2-1+j,\tilde y)$
for $0\le k\le n-1$, and
\[
\ip{h_V}w=\sum_{\tilde y}\sum_{k=0}^{n-1}W(x_2+k,\tilde y)
=-\sum_{\tilde y}\sum_{j=0}^{n-1}(n-j)D_j(\tilde y),
\]
with $D_j(\tilde y)=D(x_2-1+j,\tilde y)$ and
the sums over $\tilde y\in I_3\times\dots\times I_d$.
The value $(d_1w)(y;1,2)$ is nonzero only if one of the links $(y,1)$, $(y,2)$,
$(y+e_1,2)$, $(y+e_2,1)$ lies in $E_x$, which forces $y_1\in I_1\cup\{x_1-1\}$,
a set of $n+1$ residues; so $D(t,\tilde y)^2\le(n+1)\sum_{y_1}(d_1w)(y_1,t,
\tilde y;1,2)^2$ by the Cauchy--Schwarz inequality, and the plaquettes
$(y;1,2)$ with $y_2=x_2-1+j$, $0\le j\le n-1$, are distinct. By the
Cauchy--Schwarz inequality again,
\begin{align*}
\ip{h_V}w^2&\le n^{d-2}\sum_{j=0}^{n-1}(n-j)^2\cdot(n+1)\lVert d_1w\rVert^2\\
&=\frac{n^{d-1}(n+1)^2(2n+1)}6\lVert d_1w\rVert^2\le n^{d+2}\lVert d_1w
\rVert^2,
\end{align*}
since $6n^3-(n+1)^2(2n+1)=(n-2)(4n^2+3n+2)+3>0$ for $n\ge2$. This
is~\eqref{eq:ffclaim}. Now, with $w$ ranging over the fields supported in
$E_x$ and $t=\ip{h_V}w$, so that $\lVert w\rVert^2\ge t^2/n^d$,
\begin{align*}
\ip{h_V}{(Q_{\rm P})_{VV}^{-1}h_V}&=\sup_w\bigl[2\ip{h_V}w-\lVert d_1w\rVert^2
-\mu^2\lVert w\rVert^2\bigr]\\
&\le\sup_{t\in\R}\bigl[2t-t^2(n^{-d-2}+\mu^2n^{-d})\bigr]\\
&=\frac{n^d}{n^{-2}+\mu^2}.
\end{align*}
By Lemma~\ref{lem:fflinear}, $\ip{h_V}{(Q_{\rm P})_{VV}^{-1}h_V}$ is
$\lVert\ell_h-\mathcal P_x\ell_h\rVert^2$, the same for every $x$ by
translation invariance, and summing over the $L^d$ cubes,
\begin{align*}
\frac{n^{-d}\sum_x\lVert\ell_h-\mathcal P_x\ell_h\rVert^2}{\Var(\ell_h)}&=
\frac{\mu^2}{n^d}\ip{h_V}{(Q_{\rm P})_{VV}^{-1}h_V}\\
&\le\frac{\mu^2}{n^{-2}+\mu^2}=\frac{s^2}{1+s^2}.
\end{align*}
 Since $\ell_h$
is gauge-invariant, this bounds the infimum over gauge-invariant functions,
and the infimum over all functions is smaller.
\end{proof}

\subsubsection{Uniformity in the lattice spacing, and the size of the constants}

\begin{corollary}[the hypothesis of Theorem~\ref{thm:avgtransfer} for free
fields]\label{cor:ffuniform}
Let $m>0$ and $\ell>0$, and for a sequence of spacings $a_k\to0$ let $\mu_k=
ma_k$ and $n_k=\ell/a_k$ (or any integers $n_k\ge2$ with $\ell/2\le n_ka_k
\le\ell$). Then for every $k$ and every torus side $L\ge n_k+1$, the free
scalar field with precision $-\Delta+\mu_k^2$, and the free Proca field with
precision $\delta_1d_1+\mu_k^2$ on gauge-invariant functions, satisfy
$\mathrm{BP}(n_k,\varpi)$ with
\[
\varpi=\Bigl(\frac56\Bigr)^d\frac{(m\ell)^2}{(m\ell)^2+10d}\quad\text{if }
n_ka_k=\ell\text{ for all }k,
\]
and with
\[
\varpi=\Bigl(\frac56\Bigr)^d\frac{(m\ell)^2}{(m\ell)^2+40d}\quad\text{if only }
\ell/2\le n_ka_k\le\ell,
\]
independent of $k$ and $L$; the Proca field on all functions satisfies it
with $10d$ replaced by $5(5d-4)$ and $40d$ by $20(5d-4)$. This is the
hypothesis of Theorem~\ref{thm:avgtransfer}, with $\mathcal L_k$ the even
$L\ge2n_k+4$, for these Gaussian stand-ins of the Wilson measure. For $d=4$
and $n_ka_k=\ell$ the gauge-invariant constant is $(5/6)^4(m\ell)^2/((m\ell)^2
+40)$; by Proposition~\ref{prop:ffupper} no constant larger than
$(m\ell)^2/(1+(m\ell)^2)$ is possible for the Proca field, nor larger than
$(m\ell)^2/((m\ell)^2+16d/9)$ for the scalar field.
\end{corollary}

\begin{proof}
With $s_k=n_k\mu_k=m\,n_ka_k$, Theorems~\ref{thm:ffscalar}
and~\ref{thm:ffproca} give $c_d(s_k)$, which is increasing in $s_k$; for
$n_ka_k=\ell$ it is $c_d(m\ell)$, and for $s_k\ge m\ell/2$ it is at least
$(5/6)^d(m\ell)^2/((m\ell)^2+40d)$. The upper bound is
Proposition~\ref{prop:ffupper}(ii).
\end{proof}

\begin{remark}[the size of the constants]\label{rem:ffnumbers}
(a) In the Gaussian setting the constant of Theorem~\ref{thm:ffscalar} is
close to optimal. On the finite tori where the smallest eigenvalue of
Lemma~\ref{lem:fflinear} can be computed exactly, the $n$-dependent bound
$\alpha_n^{-d}s^2/(s^2+10d\rho_n)$ is attained at $n=2$ for every $d$ and
$s$, and lies within $2.5\%$ of the exact constant for $n=3$, $d\le4$ (and
$n=4$, $d=1$); the gauge-invariant Proca constant exceeds the bound of
Theorem~\ref{thm:ffproca}(ii) by a factor between $1.15$ and $2.0$ in
$d=2,3,4$, and the full Proca constant exceeds the bound of (i) by a factor
between $1.3$ and $3.0$; the two Proca constants coincide except at large
$s$ (for $d=4$, $n=2$, $s=4$ they are $0.5534$ and $0.5571$), the infimum
being attained at zero momentum. For the gauge-invariant Proca field in $d=4$ the
exact constant on the tori $L=4,5$ with $n=2$ is $0.0759$ at $s=1$ and
$0.2436$ at $s=2$; the bounds are $0.0588$ and $0.2$, the clean bounds
$c_4(1)=0.0118$ and $c_4(2)=0.0438$.

(b) The parabolic profile is the trial function of the proof of
Proposition~\ref{prop:gauss}(ii), and the bound of Theorem~\ref{thm:ffscalar}
is, term by term, the bound $(s_1^2/(ns_2))^d\,t/(t+2ds_1/s_2)$ proved there
for a single cube: $s_1^2/(ns_2)=1/\alpha_n$ and $2ds_1n^2/s_2=10d\rho_n$. So
the dynamics that resamples all the cubes has, on the whole of $L^2$, the
gap that Proposition~\ref{prop:gauss}(ii) gives for the Kantorovich response
of one cube; Proposition~\ref{prop:avgsup}(i) is the statement, for the
Wilson measure, that a one-cube bound of this kind implies $\mathrm{BP}$, and
for the Gaussian field Theorem~\ref{thm:ffchaos} gives the conclusion
directly. The quantity $1-\alpha_n(\mu^2)$ of Proposition~\ref{prop:gauss} is
$\mu^2n^{-d}\ip{\mathbf 1_V}{(-\Delta_D+\mu^2)^{-1}\mathbf 1_V}$, the upper
bound of Proposition~\ref{prop:ffupper}(i): it is the value of the quotient
in~\eqref{eq:ffvarpi} at the constant linear function $\ell_{\mathbf 1}$,
and in every case computed ($d\le4$, $n\le4$, $L\le16$, and $L\to\infty$ on
a momentum grid) it coincides with the exact constant to all digits, that is,
the infimum is attained at $h=\mathbf 1$. This is not proved here; what is
proved is that the two bounds squeeze $\varpi$ to within the gap between the
parabolic trial function and the exact Dirichlet resolvent, about $4\%$ for
$n\le4$.

(c) Corollary~\ref{cor:ffuniform} shows that the hypothesis of
Theorem~\ref{thm:avgtransfer}, the uniform condition $\mathrm{BP}(\ell/a_k,
\varpi)$ along a continuum tuning, is satisfied by the free massive fields
with $\varpi$ of order $(m\ell)^2$ for $m\ell\le1$; a mass gap therefore
implies the cube condition for these fields, with an explicit constant, and
the constant tends to $0$ with $m\ell$ by Proposition~\ref{prop:ffupper}; the
test function there is the zero mode, which has no Gaussian measure in the
massless case.
Inserted into the conclusion $m_*=16\varpi/(625e^2\ell)$ of
Theorem~\ref{thm:avgtransfer}, the value $\varpi=c_4(1)=0.0118$ at $m\ell=1$
gives $m_*\approx4.1\times10^{-5}\,m$: about two orders of magnitude are
the size of $\varpi$ itself, and the factor $16/(625e^2)\approx3.5\times
10^{-3}$ comes from the finite-speed argument of Lemma~\ref{lem:avgspeed}.
\end{remark}
